% Compilar con pdfLaTeX (dos pasadas para resolver las citas).
\documentclass[10pt,letterpaper]{article}
\usepackage[T1]{fontenc}
\usepackage{lmodern}
\usepackage[utf8]{inputenc}
\IfFileExists{spanish.ldf}{\usepackage[spanish]{babel}}{}
\usepackage[left=2.65cm,right=2.2cm,top=1.55cm,bottom=2.2cm]{geometry}
\usepackage{graphicx,xcolor}
\usepackage[hidelinks]{hyperref}
\setlength{\parindent}{0pt}
\setlength{\parskip}{5pt}
\pagestyle{plain}
\newcommand{\campo}[1]{\textcolor[gray]{0.36}{\textit{[#1]}}}

% Datos de la entrega: editar aquí.
\newcommand{\universidad}{Pontificia Universidad Católica de Chile}
\newcommand{\escuela}{Escuela de Ingeniería}
\newcommand{\curso}{IIC-3694 Inteligencia Artificial Generativa}
\newcommand{\titulo}{\campo{Título del proyecto}}
\newcommand{\entrega}{Entrega 1 --- Propuesta inicial}
\newcommand{\autores}{\campo{Nombre y apellido de cada integrante}}
\newcommand{\fecha}{\campo{día de mes de año}}

\begin{document}
\thispagestyle{empty}
\begin{minipage}[c]{1.95cm}
  \centering\includegraphics[height=2.4cm]{escudo-uc.png}
\end{minipage}\hspace{0.35cm}%
\begin{minipage}[c]{\dimexpr\linewidth-2.4cm\relax}
  \universidad\\
  \escuela\\
  \curso
\end{minipage}

\begin{center}
  {\large \titulo\par}
  \hrule
  \vspace{5pt}
  \entrega\par
  \autores\enspace$\cdot$\enspace\fecha\par
  \vspace{3pt}
  \hrule
\end{center}

\section{Introducción y pregunta de investigación}
\campo{Describe el problema y por qué vale la pena estudiarlo. Formula una pregunta que puedas responder con los experimentos del proyecto.}

\section{Revisión bibliográfica}
\campo{Explica qué problema aborda el primer trabajo}~\cite{trabajo1}.
\campo{Contrasta su método o sus resultados con los del segundo}~\cite{trabajo2}.
\campo{Identifica qué aspecto queda abierto y cómo orienta tu propuesta.}

\section{Propuesta de investigación}
\campo{Plantea una hipótesis tentativa, el objetivo del estudio y el resultado que permitiría apoyar o cuestionar esa hipótesis.}

\section{Entorno y diseño experimental}
\campo{Describe la tarea, el modelo base, los datos, las particiones y los niveles de dificultad. Indica cómo el verificador determina si una respuesta es correcta.}

\campo{Explica cómo compararás el modelo base, solo SFT, SFT seguido de RL, solo RL y DPO con un presupuesto de cómputo comparable. Define el algoritmo de RL, las semillas y el valor de k para reportar pass@1 y pass@k por dificultad.}

% Para la entrega 2, cambia el campo de entrega y descomenta estas secciones.
% \section{Implementación inicial}
% \campo{Describe qué componentes están implementados y cómo reproducir la evaluación.}
% \section{Resultados preliminares}
% \campo{Reporta la línea base antes de entrenar. Inserta las primeras curvas de recompensa e interpreta si el entorno entrega señal de aprendizaje.}
% \section{Iteraciones y próximos pasos}
% \campo{Documenta un problema observado en las trazas, la corrección y su efecto.}

% Para el paper final, trasladar el contenido al formato que publique el curso.
% Incluir diagrama de bloques, resultados con el verificador final, auditoría
% manual de trazas exitosas, ciclo de iteración y conclusiones.
\renewcommand{\refname}{Referencias}
\begin{thebibliography}{9}
  \bibitem{trabajo1} \campo{Autores}. \campo{Título del primer trabajo}. \campo{Revista o conferencia}, \campo{año}.
  \bibitem{trabajo2} \campo{Autores}. \campo{Título del segundo trabajo}. \campo{Revista o conferencia}, \campo{año}.
\end{thebibliography}
\end{document}
